% !TEX program = xelatex
% 系泊系统结构及主要几何符号示意图（非比例示意）。
% 编译：xelatex mooring_symbols.tex
\documentclass[tikz,border=5mm]{standalone}
\usepackage[UTF8,fontset=fandol]{ctex}
\usepackage{amsmath,bm}
\usepackage{tikz}
\usetikzlibrary{arrows.meta,calc,angles,quotes,positioning}

\begin{document}
\begin{tikzpicture}[
  x=0.75cm,y=0.75cm,
  >=Stealth,
  font=\small,
  water/.style={black!70,line width=0.45pt},
  bed/.style={black,line width=0.55pt},
  chain/.style={black!85,line width=0.8pt,line cap=round},
  component/.style={black,line width=5.2pt,line cap=round},
  pipe/.style={black,line width=1.15pt,line cap=round},
  joint/.style={circle,fill=black,inner sep=1.25pt},
  dim/.style={black!75,thin,<->},
  extension/.style={black!55,densely dashed,thin},
  reference/.style={black!55,densely dashed,thin},
  anglemark/.style={draw=black!80,thin,angle radius=0.30cm,angle eccentricity=1.55},
  labeltxt/.style={font=\small,inner sep=1.2pt,fill=white,fill opacity=0.88,text opacity=1},
  noticetxt/.style={font=\scriptsize,inner sep=1pt,fill=white,fill opacity=0.88,text opacity=1}
]

% --------------------------------------------------------------------------
% 环境与样式：下列命名参数只控制示意图版式，不代表计算工况或比例。
% --------------------------------------------------------------------------
\def\SeaY{8.40}
\def\BedY{0.85}
\def\AnchorX{3.38}
\def\BuoyX{11.10}
\def\BuoyBottomY{7.72}

% --------------------------------------------------------------------------
% 海面和海床
% --------------------------------------------------------------------------
\coordinate (seaLeft) at (0.55,\SeaY);
\coordinate (seaRight) at (15.70,\SeaY);
\coordinate (bedLeft) at (0.55,\BedY);
\coordinate (bedRight) at (15.70,\BedY);
\draw[water] (seaLeft) -- (seaRight);
\draw[bed] (bedLeft) -- (bedRight);

% --------------------------------------------------------------------------
% 锚与锚链：实线链形仅为平滑结构示意，不代表精确悬链线计算结果。
% --------------------------------------------------------------------------
\coordinate (anchorSW) at (2.58,\BedY);
\coordinate (anchorNE) at (3.38,1.37);
\coordinate (anchorPoint) at (3.38,1.05);
\coordinate (touchdown) at (5.55,1.05);
\coordinate (chainControlA) at (7.05,1.12);
\coordinate (chainControlB) at (8.85,2.18);
\coordinate (chainJoint) at (10.12,3.83);
\coordinate (chainLabel) at (7.62,2.62);
\draw[fill=black!88,draw=black] (anchorSW) rectangle (anchorNE);
\draw[chain] (anchorPoint) -- (touchdown)
  .. controls (chainControlA) and (chainControlB) .. (chainJoint);

% --------------------------------------------------------------------------
% 钢桶和重物球：锚链、钢桶和悬挂支线共用 chainJoint 连接点。
% --------------------------------------------------------------------------
\coordinate (bucketLower) at (chainJoint);
\coordinate (bucketUpper) at (10.54,4.86);
\coordinate (bucketMid) at ($(bucketLower)!0.5!(bucketUpper)$);
\coordinate (bucketVertical) at ($(bucketMid)+(0,0.92)$);
\coordinate (weightTop) at (10.12,3.12);
\coordinate (weightCenter) at (10.12,2.74);
\draw[component] (bucketLower) -- (bucketUpper);
\draw[chain] (chainJoint) -- (weightTop);
\fill[black] (weightCenter) circle[radius=0.28];

% --------------------------------------------------------------------------
% 四节钢管：从浮标向钢桶编号 i=1,2,3,4。
% --------------------------------------------------------------------------
\coordinate (pipeFourLower) at (bucketUpper);
\coordinate (jointFourThree) at (10.70,5.62);
\coordinate (jointThreeTwo) at (10.88,6.36);
\coordinate (jointTwoOne) at (11.00,7.06);
\coordinate (pipeOneUpper) at (\BuoyX,\BuoyBottomY);
\coordinate (pipeOneMid) at ($(jointTwoOne)!0.5!(pipeOneUpper)$);
\coordinate (pipeTwoMid) at ($(jointThreeTwo)!0.5!(jointTwoOne)$);
\coordinate (pipeThreeMid) at ($(jointFourThree)!0.5!(jointThreeTwo)$);
\coordinate (pipeFourMid) at ($(pipeFourLower)!0.5!(jointFourThree)$);
\coordinate (pipeOneVertical) at ($(pipeOneMid)+(0,0.58)$);
\coordinate (pipeTwoVertical) at ($(pipeTwoMid)+(0,0.58)$);
\coordinate (pipeThreeVertical) at ($(pipeThreeMid)+(0,0.58)$);
\coordinate (pipeFourVertical) at ($(pipeFourMid)+(0,0.58)$);
\coordinate (thetaOneLabel) at ($(pipeOneMid)+(-0.88,-0.08)$);
\coordinate (thetaTwoLabel) at ($(pipeTwoMid)+(-0.88,-0.04)$);
\coordinate (thetaThreeLabel) at ($(pipeThreeMid)+(-0.88,-0.04)$);
\coordinate (thetaFourLabel) at ($(pipeFourMid)+(-0.88,-0.04)$);
\draw[pipe] (pipeFourLower) -- (jointFourThree);
\draw[pipe] (jointFourThree) -- (jointThreeTwo);
\draw[pipe] (jointThreeTwo) -- (jointTwoOne);
\draw[pipe] (jointTwoOne) -- (pipeOneUpper);
\node[joint] at (jointFourThree) {};
\node[joint] at (jointThreeTwo) {};
\node[joint] at (jointTwoOne) {};

% --------------------------------------------------------------------------
% 浮标
% --------------------------------------------------------------------------
\coordinate (buoySW) at (10.43,\BuoyBottomY);
\coordinate (buoyNE) at (11.77,9.55);
\coordinate (buoyAxisBottom) at (\BuoyX,\BuoyBottomY);
\coordinate (buoyAxisSea) at (\BuoyX,\SeaY);
\coordinate (buoyAxisTop) at (\BuoyX,9.55);
\path[fill=black!12,draw=black,line width=0.65pt,rounded corners=2.5pt]
  (buoySW) rectangle (buoyNE);

% --------------------------------------------------------------------------
% 尺寸标注：H、delta、r、L 和 L_0。
% --------------------------------------------------------------------------
\coordinate (depthSea) at (1.02,\SeaY);
\coordinate (depthBed) at (1.02,\BedY);
\coordinate (depthSeaExt) at (0.55,\SeaY);
\coordinate (depthBedExt) at (0.55,\BedY);
\draw[extension] (depthSea) -- (depthSeaExt);
\draw[extension] (depthBed) -- (depthBedExt);
\draw[dim] (depthBed) -- (depthSea) node[midway,labeltxt,left] {$H$};

\coordinate (draftSea) at (12.25,\SeaY);
\coordinate (draftBottom) at (12.25,\BuoyBottomY);
\coordinate (draftSeaExt) at (11.77,\SeaY);
\coordinate (draftBottomExt) at (11.77,\BuoyBottomY);
\draw[extension] (draftSea) -- (draftSeaExt);
\draw[extension] (draftBottom) -- (draftBottomExt);
\draw[dim] (draftBottom) -- (draftSea) node[midway,labeltxt,right] {$\delta$};

\coordinate (rangeLeft) at (\AnchorX,0.34);
\coordinate (rangeRight) at (\BuoyX,0.34);
\coordinate (rangeLeftExt) at (\AnchorX,\BedY);
\coordinate (rangeRightExt) at (\BuoyX,\BuoyBottomY);
\draw[extension] (rangeLeft) -- (rangeLeftExt);
\draw[extension] (rangeRight) -- (rangeRightExt);
\draw[dim] (rangeLeft) -- (rangeRight) node[midway,noticetxt,below] {$r$};

\coordinate (bedChainDimLeft) at ($(anchorPoint)+(0,0.47)$);
\coordinate (bedChainDimRight) at ($(touchdown)+(0,0.47)$);
\draw[dim] (bedChainDimLeft) -- (bedChainDimRight)
  node[midway,noticetxt,above] {$L_0$};
\coordinate (groundedConditionLabel) at (4.52,1.96);
\node[labeltxt] at (chainLabel) {锚链总弧长 $L$};
\node[noticetxt] at (groundedConditionLabel)
  {$L_0>0,\ \varphi_{\mathrm a}=0$};

% --------------------------------------------------------------------------
% 倾角标注：钢管与钢桶轴线均相对局部竖直参考线定义。
% --------------------------------------------------------------------------
\draw[reference] (pipeOneMid) -- (pipeOneVertical);
\draw[reference] (pipeTwoMid) -- (pipeTwoVertical);
\draw[reference] (pipeThreeMid) -- (pipeThreeVertical);
\draw[reference] (pipeFourMid) -- (pipeFourVertical);
\draw[reference] (bucketMid) -- (bucketVertical);
\pic[anglemark] {angle=pipeOneUpper--pipeOneMid--pipeOneVertical};
\pic[anglemark] {angle=jointTwoOne--pipeTwoMid--pipeTwoVertical};
\pic[anglemark] {angle=jointThreeTwo--pipeThreeMid--pipeThreeVertical};
\pic[anglemark] {angle=jointFourThree--pipeFourMid--pipeFourVertical};
\pic[anglemark] {angle=bucketUpper--bucketMid--bucketVertical};
\node[noticetxt,anchor=west] at (thetaOneLabel) {$\theta_1$};
\node[noticetxt,anchor=west] at (thetaTwoLabel) {$\theta_2$};
\node[noticetxt,anchor=west] at (thetaThreeLabel) {$\theta_3$};
\node[noticetxt,anchor=west] at (thetaFourLabel) {$\theta_4$};
\node[noticetxt,anchor=south west] at ($(bucketMid)+(0.14,0.26)$) {$\theta_{\mathrm d}$};

% --------------------------------------------------------------------------
% 构件名称与环境参数
% --------------------------------------------------------------------------
\coordinate (windStart) at (3.55,9.18);
\coordinate (windEnd) at (5.80,9.18);
\coordinate (currentStart) at (3.70,6.30);
\coordinate (currentEnd) at (5.75,6.30);
\coordinate (rhoLabel) at (5.05,5.30);
\draw[->,black!75,thin] (windStart) -- (windEnd)
  node[midway,noticetxt,above] {$v_{\mathrm w}$};
\draw[->,black!75,thin] (currentStart) -- (currentEnd)
  node[midway,noticetxt,above] {$v_{\mathrm c}$};
\node[labeltxt] at (rhoLabel) {海水密度 $\rho$};
\node[labeltxt,anchor=west] at (1.22,8.72) {海面};
\node[labeltxt,anchor=west] at (12.70,0.85) {海床};
\node[labeltxt,anchor=east] at ($(anchorSW)+(-0.18,0.28)$) {锚};
\node[labeltxt,anchor=west] at ($(bucketMid)+(0.58,0.05)$) {钢桶};
\node[labeltxt,anchor=west] at ($(weightCenter)+(0.48,0.00)$) {重物球 $m_{\mathrm s}$};
\node[labeltxt,anchor=west] at ($(pipeThreeMid)+(-1.56,0.18)$) {钢管};
\node[labeltxt,anchor=west] at ($(buoyNE)+(0.14,-0.18)$) {浮标};
\node[noticetxt,anchor=west] at ($(pipeOneMid)+(0.62,0.20)$) {$i=1$};
\node[noticetxt,anchor=west] at ($(pipeTwoMid)+(0.62,0.05)$) {$i=2$};
\node[noticetxt,anchor=west] at ($(pipeThreeMid)+(0.62,-0.08)$) {$i=3$};
\node[noticetxt,anchor=west] at ($(pipeFourMid)+(0.60,-0.16)$) {$i=4$};

% --------------------------------------------------------------------------
% 锚端局部插图：仅定义完全离床的 L_0=0, varphi_a>0 情形。
% --------------------------------------------------------------------------
\coordinate (insetBedLeft) at (12.78,1.28);
\coordinate (insetBedRight) at (15.72,1.28);
\coordinate (insetAnchorSW) at (13.02,1.28);
\coordinate (insetAnchorNE) at (13.28,1.50);
\coordinate (insetAnchorPoint) at (13.28,1.43);
\coordinate (insetHorizontal) at (14.02,1.43);
\coordinate (insetTangent) at (13.96,1.86);
\coordinate (insetCurveEnd) at (15.12,2.52);
\draw[black!65,thin] (insetBedLeft) -- (insetBedRight);
\draw[fill=black!88,draw=black] (insetAnchorSW) rectangle (insetAnchorNE);
\draw[chain] (insetAnchorPoint) -- (insetTangent)
  .. controls (14.28,2.05) and (14.72,2.36) .. (insetCurveEnd);
\pic[anglemark,"$\varphi_{\mathrm a}$" {noticetxt,anchor=south west}]
  {angle=insetHorizontal--insetAnchorPoint--insetTangent};
\node[noticetxt,anchor=west] at (12.78,2.94) {锚端局部状态（完全离床）};
\node[noticetxt,anchor=west] at (12.78,2.62) {$L_0=0,\ \varphi_{\mathrm a}>0$};

\node[noticetxt] at (7.10,-0.20) {示意图，未按比例绘制};
\end{tikzpicture}
\end{document}
